% Figura: problema 4.3. Condensador plano con dieléctrico: cargas libres, ligadas,
% campos y superficies gaussianas.
\begin{figure}[t]
\centering
\begin{tikzpicture}[scale=1.0]
  % dieléctrico y placas
  \fill[dielec] (0,0) rectangle (7.0,2.6);
  \fill[conductor] (0,2.6) rectangle (7.0,2.82);
  \fill[conductor] (0,-0.22) rectangle (7.0,0);
  \draw[texto, line width=0.5pt] (0,2.6) rectangle (7.0,2.82) (0,-0.22) rectangle (7.0,0);
  \foreach \x in {0.3,0.85,...,6.9} {
    \node[cneg, font=\scriptsize\bfseries] at (\x,2.71) {$-$};
    \node[cpos, font=\scriptsize\bfseries] at (\x,-0.11) {$+$};
    \node[cpos, font=\tiny\bfseries] at (\x,2.47) {$+$};
    \node[cneg, font=\tiny\bfseries] at (\x,0.13) {$-$};
  }
  \node[etiq, anchor=west] at (7.1,2.71) {$-\sigma$};
  \node[etiq, anchor=west] at (7.1,-0.11) {$+\sigma$};
  \node[etiq, anchor=west, text=acentoC] at (7.1,2.42) {$+\sigmaP$};
  \node[etiq, anchor=west, text=acentoC] at (7.1,0.18) {$-\sigmaP$};
  % campos
  \draw[vecD] (0.65,0.35) -- (0.65,2.2) node[pos=0.5, left, etiq] {$\vD$};
  \draw[vecE] (1.35,0.84) -- (1.35,1.76) node[pos=0.5, left, etiq] {$\vE$};
  \draw[vecP] (2.0,0.84) -- (2.0,1.76) node[pos=0.5, left, etiq] {$\vP$};
  \node[etiq] at (2.85,1.3) {$\varepsilon$};
  % superficie gaussiana S (encierra la placa superior)
  \draw[gauss] (3.5,1.5) rectangle (4.7,3.2);
  \node[etiq, text=acentoB, anchor=west] at (4.75,3.05) {$S$};
  \draw[vecn] (4.1,1.5) -- (4.1,0.95) node[right, etiq] {$\nn$};
  % superficie gaussiana S' (solo la cara del dieléctrico y la placa inferior)
  \draw[gauss] (5.3,-0.5) rectangle (6.5,1.05);
  \node[etiq, text=acentoB, anchor=west] at (6.55,0.95) {$S'$};
  \draw[vecn] (5.9,1.05) -- (5.9,1.6) node[right, etiq] {$\nn'$};
  % cotas
  \draw[cota] (-0.4,0) -- (-0.4,2.6) node[midway, left, etiq] {$d$};
  \node[etiq, anchor=south] at (1.6,2.86) {$A$};
\end{tikzpicture}
\caption{\Cref{prob:condensador-dielectrico}. Condensador plano con cargas libres $\pm\sigma$ en las placas y un dieléctrico de
permitividad $\varepsilon$. Junto a cada placa aparece una carga ligada de signo opuesto, $\pm\sigmaP$ con
$\sigmaP=\sigma(\varepsilon-\eo)/\varepsilon<\sigma$. Las tres flechas indican el sentido común de $\vD$, $\vE$ y $\vP$ (de la placa positiva a
la negativa); sus longitudes son proporcionales a $D=\sigma$, $\eo E=\sigma/\er$ y $P=\sigmaP$ (dibujadas con $\er=2$). Las superficies
$S$ y $S'$ sirven para aplicar la ley de Gauss para $\vD$ (solo cuenta la carga libre).}
\label{fig:p43}
\end{figure}
